\documentclass[11pt,a4paper]{article}
\usepackage[margin=2.5cm]{geometry}
\usepackage{amsmath,amssymb,booktabs,array}
\usepackage[hidelinks]{hyperref}

\newcommand{\dd}{\mathrm{d}}
\newcommand{\Msq}{|\mathcal{M}|^2}
\newcommand{\code}[1]{\texttt{#1}}

\title{\texorpdfstring{Integration algorithms for the $2\leftrightarrow2$ collision integral\\
\large \code{coll\_12\_34} and \code{coll\_12\_34\_legendre}}{Integration algorithms for the 2-2 collision integral}}
\author{}
\date{\today}

\begin{document}
\maketitle

\section{The quantity computed}

Both functions compute, for a process $1\,2\leftrightarrow3\,4$ (set up for
$A'A'\leftrightarrow NN$, i.e.\ $m_1=m_2=m_{A'}$, $m_3=m_4=m_N$),
\begin{equation}
  \mathcal{C} = \int \dd\Pi_1\,\dd\Pi_2\,\dd\Pi_3\,\dd\Pi_4\,
  (2\pi)^4\delta^4(p_1+p_2-p_3-p_4)\,\Msq\,F ,
  \qquad \dd\Pi_i = \frac{\dd^3p_i}{(2\pi)^3 2E_i},
  \label{eq:C}
\end{equation}
with
\begin{equation}
  F = f_1 f_2 (1-k_3 f_3)(1-k_4 f_4) - f_3 f_4 (1-k_1 f_1)(1-k_2 f_2),
  \qquad
  f_i = \frac{1}{e^{E_i/T_i-\xi_i}+k_i},
\end{equation}
$k_i=+1$ ($-1$) for fermions (bosons), and $\Msq$ summed over all degrees of
freedom (eq.~(A9) of the paper, including $g^4$). No symmetry factors for
identical particles are included. Throughout,
\begin{gather}
  E \equiv E_1+E_2 = E_3+E_4,\qquad
  p_i\equiv|\mathbf{p}_i| = \sqrt{E_i^2-m_i^2},\qquad
  P\equiv|\mathbf{p}_1+\mathbf{p}_2| = \sqrt{E^2-s},\\
  \lambda_{ab}\equiv\lambda(s,m_a^2,m_b^2)
  = \big(s-(m_a+m_b)^2\big)\big(s-(m_a-m_b)^2\big).
\end{gather}

\paragraph{Common reduction.}
Integrating $\mathbf{p}_4$ against $\delta^3$, writing
$\dd^3p_i = p_iE_i\,\dd E_i\,\dd\Omega_i$, trading the angle between
$\mathbf{p}_1$ and $\mathbf{p}_2$ for $s$, the polar angle of $\mathbf{p}_3$
around $\mathbf{P}$ for the energy $\delta$-function, and keeping the azimuth
$\phi$ of $\mathbf{p}_3$ around $\mathbf{P}$ (measured from the
$(\mathbf{p}_1,\mathbf{P})$ plane) gives
\begin{equation}
  \mathcal{C} = \frac{1}{512\pi^6}\int\dd E_1\,\dd E_2\,\dd E_3\;F
  \int\dd s\;\frac{1}{P}\int_0^\pi\dd\phi\;\Msq(s,t).
  \label{eq:reduced}
\end{equation}
The domain is $E_i\ge m_i$ and $s$ such that both triangles
$\mathbf{P}=\mathbf{p}_1+\mathbf{p}_2=\mathbf{p}_3+\mathbf{p}_4$ close. The
two algorithms differ in how they parametrize and integrate this domain.

%------------------------------------------------------------------------------
\section{\texorpdfstring{Algorithm 1: nested adaptive quadrature (\code{coll\_12\_34})}{Algorithm 1: nested adaptive quadrature}}

\subsection{Integration variables and order}

The integration order, from outermost to innermost, is
$E_1\to E_2\to E_3\to s\to t$, with $E_4=E_1+E_2-E_3$. Here $t=(p_1-p_3)^2$
replaces $\phi$ in \eqref{eq:reduced}. The formula implemented is
\begin{equation}
  \mathcal{C} = \frac{g^4}{256\pi^6}
  \int_{E_1^-}^{E_1^+}\!\dd E_1
  \int_{E_2^-}^{E_2^+}\!\dd E_2
  \int_{E_3^-}^{E_3^+}\!\dd E_3\;p_3\,F
  \int_{s^-}^{s^+}\!\dd s
  \int_{t^-}^{t^+}\!\dd t\;
  \frac{\Msq/g^4}{\sqrt{-a\,(t-t^-)(t^+-t)}} ,
  \label{eq:alg1}
\end{equation}
where $a<0$ is defined in Sec.~\ref{sec:alg1-limits}. This agrees with
\eqref{eq:reduced}: with $t=\bar t+h\cos\phi$ the $t$-integral equals
$(-a)^{-1/2}\int_0^\pi\dd\phi\,\Msq$, and $p_3/\sqrt{-a}=1/(2P)$.

\subsection{Limits}
\label{sec:alg1-limits}

This parametrisation (integration order $E_1,E_2,E_3,s$, then the angle
$\theta$ or equivalently $t$) and its limits are those of
Ref.~\cite{Bringmann:2022aim}; see the appendix ``Collision operators'',
Eqs.~(17)--(24) of the arXiv version. The authors there state that they
follow an approach similar to Ref.~\cite{Ala-Mattinen:2022nuj}. The code
uses the sign conventions of Ref.~\cite{Bringmann:2022aim} for the quadratic
in $\cos\theta$: the coefficients $a<0$, $b$, $c$ and the root labels
$c_{\theta,\pm}$ below are identical to theirs. Two differences remain:
\begin{itemize}
\item The inner integral runs over $t$ rather than $\cos\theta$.
\item The energy integrals end at the cutoffs given below rather than at $\infty$.
\end{itemize}
The prefactor $1/(4(2\pi)^6)=1/(256\pi^6)$ agrees with \eqref{eq:alg1}; the
symmetry factor $\kappa$ that appears there is not included in \code{coll\_12\_34}.

\begin{description}
\item[$E_1$:] $E_1^- = m_1$, $\;E_1^+ = \max(100\,T_1,\,100\,m_1)$.
  The upper limit is a cutoff.
\item[$E_2$:] $E_2^- = \max(m_2,\; m_3+m_4-E_1)$,
  $\;E_2^+ = \max(100\,T_2,\,100\,m_2)$. The upper limit is a cutoff; the
  result is $0$ if $E_2^-\ge E_2^+$.
\item[$E_3$:] $E_3^- = m_3$, $\;E_3^+ = E_1+E_2-m_4$. Both are exact;
  $E_4 = E_1+E_2-E_3$.
\item[$s$:] Both limits are exact. With
  $s_{ab}^{\pm} = (E_a+E_b)^2-(p_a\mp p_b)^2$ (code: \code{s\_lim}),
  \begin{equation}
    s^- = \max\big(s_{12}^-,\,s_{34}^-\big),\qquad
    s^+ = \min\big(s_{12}^+,\,s_{34}^+\big),
  \end{equation}
  i.e.\ $|p_1-p_2|\le P\le p_1+p_2$ and $|p_3-p_4|\le P\le p_3+p_4$. The
  result is $0$ if $s^-\ge s^+$.
\item[$t$:] Here $\theta$ is the lab angle between $\mathbf{p}_1$ and
  $\mathbf{p}_3$, so that $t = m_1^2+m_3^2-2E_1E_3+2p_1p_3\cos\theta$. The
  limits of $\cos\theta$ are the roots $c_{\theta,\pm}$ of
  $a\cos^2\theta+b\cos\theta+c=0$, with (Ref.~\cite{Bringmann:2022aim},
  Eqs.~(17)--(20))
  \begin{gather}
    a = -4p_3^2(E^2-s),\qquad
    b = 2\,\frac{p_3}{p_1}\,\delta_{12}\delta_{34},\qquad
    c = -\delta_{34}^2-\frac{p_3^2}{p_1^2}\,(s-s_{12}^-)(s-s_{12}^+),\\
    \delta_{ab} = s+m_a^2-m_b^2-2E_aE,\\
    c_{\theta,\pm} = \frac{-b\pm\sqrt{D}}{2a},\qquad
    D = b^2-4ac
    = 4\frac{p_3^2}{p_1^2}\,(s-s_{12}^-)(s-s_{12}^+)(s-s_{34}^-)(s-s_{34}^+).
  \end{gather}
  Since $a<0$, $c_{\theta,+}\le c_{\theta,-}$, and
  $a\cos^2\theta+b\cos\theta+c = -a\,(\cos\theta-c_{\theta,+})(c_{\theta,-}-\cos\theta)\ge0$
  on the integration region $c_{\theta,+}\le\cos\theta\le c_{\theta,-}$. Hence
  \begin{equation}
    t^- = m_1^2+m_3^2-2E_1E_3+2p_1p_3\,c_{\theta,+},\qquad
    t^+ = m_1^2+m_3^2-2E_1E_3+2p_1p_3\,c_{\theta,-}.
  \end{equation}
  In the code $c_{\theta,\pm}=$ \code{cos\_theta\_lim($\pm$1)}, so
  \code{t\_lim(1)} $=t^-$ and \code{t\_lim(-1)} $=t^+$. $D$ is evaluated in
  the factorized form, and $c_{\theta,\pm}$ is clamped to $[-1,1]$ against
  round-off. Geometrically, $\cos\theta_{1P}=-\delta_{12}/(2p_1P)$ and
  $\cos\theta_{3P}=-\delta_{34}/(2p_3P)$ are the fixed angles of
  $\mathbf{p}_1$ and $\mathbf{p}_3$ with $\mathbf{P}$, so that
  $c_{\theta,\pm}=\cos(\theta_{1P}\pm\theta_{3P})$.
\end{description}

\subsection{The $t$-integral (closed form)}

For $m_1=m_2$ and $m_3=m_4$ one has $t,u\le0$ throughout the physical region,
so the $N$-exchange poles at $t,u=m_N^2$ are never reached. Define the
positive quantities
\begin{equation}
  D_t^\pm = m_N^2-t^\pm,\qquad D_u^\pm = s+t^\pm-2m_{A'}^2-m_N^2 = m_N^2-u(t^\pm).
\end{equation}
Substituting $t=\bar t+h\cos\phi$ with $\phi\in[0,\pi]$, and using
$D_t(\phi)+D_u(\phi)=s-2m_{A'}^2$ together with
$\int_0^\pi\dd\phi/(\alpha+\beta\cos\phi)=\pi/\sqrt{\alpha^2-\beta^2}$ and its
derivative with respect to $\alpha$:
\begin{multline}
  \int_{t^-}^{t^+}\!\!\frac{\dd t\;\Msq/g^4}{\sqrt{-a(t-t^-)(t^+-t)}}
  = \frac{8\pi}{\sqrt{-a}}\Bigg[
    \frac{N_{tu}}{s-2m_{A'}^2}\Big(\frac{1}{\sqrt{D_t^+D_t^-}}+\frac{1}{\sqrt{D_u^+D_u^-}}\Big)\\
    - C^2\,\frac{(D_t^++D_t^-)/2}{(D_t^+D_t^-)^{3/2}}
    - C^2\,\frac{(D_u^++D_u^-)/2}{(D_u^+D_u^-)^{3/2}}
    - 2\Bigg],
\end{multline}
with $N_{tu}=4(m_{A'}^2-m_N^2)^2-16m_N^4+(2m_N^2+s)^2$ and
$C=m_{A'}^2+2m_N^2$.

\subsection{Quadrature}

Each of the four remaining integrals uses adaptive Gauss--Kronrod quadrature
with 7/15 points (\code{QuadGKJL} from \code{Integrals.jl}) on the
untransformed variable, with
user-set \code{reltol} (default $10^{-6}$) and \code{abstol} (default $0$).
Only interior nodes are used, so the vanishing of $a$ at $E_3=m_3$ ($p_3=0$)
does no harm: the combination $p_3/\sqrt{-a}$ stays finite. The $s$-integrand is
skipped (set to 0) if both terms of $F$ are $<10^{-30}$, or if they cancel to
relative precision $10^{-13}$.

%------------------------------------------------------------------------------
\section{\texorpdfstring{Algorithm 2: CM-frame Legendre decomposition\\ (\code{coll\_12\_34\_legendre})}{Algorithm 2: CM-frame Legendre decomposition}}

\subsection{Integration variables}

Fix $s$ and $E$, and go to the CM frame, which moves with velocity
$\beta=P/E$ ($\gamma=E/\sqrt s$) along $\hat{\mathbf{P}}$. CM quantities are
\begin{gather}
  E_1^*=\frac{s+m_1^2-m_2^2}{2\sqrt s},\qquad
  E_3^*=\frac{s+m_3^2-m_4^2}{2\sqrt s},\qquad
  E_4^*=\sqrt s-E_3^*,\\
  k_1=|\mathbf{p}_1^*|=|\mathbf{p}_2^*|=\frac{\lambda_{12}^{1/2}}{2\sqrt s},\qquad
  k_3=|\mathbf{p}_3^*|=|\mathbf{p}_4^*|=\frac{\lambda_{34}^{1/2}}{2\sqrt s}.
\end{gather}
The variables are:
\begin{itemize}
\item $c_1=\cos\angle(\mathbf{p}_1^*,\hat{\mathbf{P}})$ and
  $c_3=\cos\angle(\mathbf{p}_3^*,\hat{\mathbf{P}})$, both in $[-1,1]$. The
  lab energies are linear in them:
  \begin{equation}
    E_1 = \gamma(E_1^*+\beta k_1c_1) = \frac{E(s+m_1^2-m_2^2)+c_1P\lambda_{12}^{1/2}}{2s},
    \qquad
    E_3 = \frac{E(s+m_3^2-m_4^2)+c_3P\lambda_{34}^{1/2}}{2s},
  \end{equation}
  with $E_2=E-E_1$ and $E_4=E-E_3$. In the code each $E_i$ is clamped to
  $\ge m_i$ against round-off.
\item $\theta^*\in[0,\pi]$, $c^*=\cos\theta^*=\hat{\mathbf{p}}_1^*\cdot\hat{\mathbf{p}}_3^*$,
  the CM scattering angle. In terms of it,
  \begin{gather}
    t = t_0-4k_1k_3\sin^2(\theta^*/2),\qquad u = u_0-4k_1k_3\cos^2(\theta^*/2),\\
    t_0 = m_1^2+m_3^2-\frac{2(m_1^2E_3^{*2}+m_3^2k_1^2)}{E_1^*E_3^*+k_1k_3},\qquad
    u_0 = m_1^2+m_4^2-\frac{2(m_1^2E_4^{*2}+m_4^2k_1^2)}{E_1^*E_4^*+k_1k_3},
  \end{gather}
  which avoids cancellation near the forward and backward peaks.
\item $w=\sqrt s = w_{\rm th}+\sigma^2$, with $w_{\rm th}=\max(m_1+m_2,\,m_3+m_4)$
  and $\sigma\ge0$. This makes the threshold factor $\lambda^{1/2}\propto\sigma$
  smooth; $\dd s = 4w\sigma\,\dd\sigma$.
\item $E = w+\tau^2$ with $\tau\ge0$. Then $P=\tau\sqrt{\tau^2+2w}$ and
  $P\,\dd E = 2\tau P\,\dd\tau$, which removes the $\sqrt{E-\sqrt s}$ endpoint.
\end{itemize}
At fixed $(s,E)$, $\dd E_1\,\dd E_3 = P^2\lambda_{12}^{1/2}\lambda_{34}^{1/2}/(4s^2)\;\dd c_1\,\dd c_3$,
and $(c_3,\phi)$ are spherical coordinates of $\hat{\mathbf{p}}_3^*$ about
$\hat{\mathbf P}$, since boosts along $\hat{\mathbf P}$ preserve the azimuth.
Equation \eqref{eq:reduced} becomes
\begin{gather}
  \mathcal{C} = \frac{1}{512\pi^6}\int\dd s\int\dd E\;
  \frac{P\,\lambda_{12}^{1/2}\lambda_{34}^{1/2}}{4s^2}
  \int_{-1}^1\!\dd c_1\int_{-1}^1\!\dd c_3\int_0^\pi\!\dd\phi\;\Msq(s,c^*)\,F,\\
  c^*=c_1c_3+\sqrt{1-c_1^2}\sqrt{1-c_3^2}\cos\phi .
\end{gather}

\subsection{Angular integrals via Funk--Hecke}

At fixed $(s,E)$ the distribution factor splits into products,
$F=a(c_1)\,b(c_3)-a'(c_1)\,b'(c_3)$, with
\begin{equation}
  a = f_1f_2,\quad b=(1-k_3f_3)(1-k_4f_4),\quad
  a'=(1-k_1f_1)(1-k_2f_2),\quad b' = f_3f_4 .
\end{equation}
Expanding in Legendre polynomials and applying the Funk--Hecke theorem (proven e.g.\ in
Ref.~\cite{Mueller:1998}) twice gives
\begin{gather}
  \int_{-1}^1\!\dd c_1\int_{-1}^1\!\dd c_3\int_0^\pi\!\dd\phi\;\Msq\,a\,b
  = \frac{\pi}{2}\sum_{l=0}^\infty(2l+1)\,M_l(s)\,\hat a_l\,\hat b_l,\\
  M_l(s) = \int_{-1}^1\dd c^*\,\Msq(s,c^*)\,P_l(c^*),\qquad
  \hat a_l = \int_{-1}^1\dd c\;a(c)\,P_l(c),\qquad
  \hat b_l = \int_{-1}^1\dd c\;b(c)\,P_l(c).
\end{gather}
Expanding the angular dependence of a $2\leftrightarrow2$ collision kernel
in Legendre polynomials and truncating the series is an established technique
in neutrino transport. Ref.~\cite{Pons:1998} does this for the pair-annihilation
kernel $\nu\bar\nu\leftrightarrow e^+e^-$ and studies how the result depends
on the truncation order. There the Legendre moments are those of anisotropic
neutrino distributions. Here the distributions are isotropic in the lab frame,
and the angular structure instead comes from the dependence of the lab
energies on the CM angles $c_1$, $c_3$.
The result is
\begin{equation}
  \boxed{\;
  \mathcal{C} = \frac{1}{4096\pi^5}\int_0^{\sigma_{\max}}\!\dd\sigma\;
  4w\sigma\,\frac{\lambda_{12}^{1/2}\lambda_{34}^{1/2}}{s^2}
  \sum_{l=0}^{L}(2l+1)\,M_l(s)
  \int_0^{\tau_{\max}}\!\dd\tau\;2\tau P\,
  \big[\hat a_l\hat b_l-\hat a'_l\hat b'_l\big]\;}
  \label{eq:alg2}
\end{equation}
In the Maxwell--Boltzmann limit $a,b$ are constant, only $l=0$ survives, and
\eqref{eq:alg2} reduces to
$\frac{T}{512\pi^5}\int\dd s\,K_1(\sqrt s/T)\,s^{-1/2}\int\dd t\,\Msq$,
up to the chemical-potential factors.

\paragraph{Equal temperatures.}
If $T_1=\dots=T_4$, energy conservation gives $a'b'=e^{\Delta\xi}\,ab$
pointwise, with $\Delta\xi=\xi_3+\xi_4-\xi_1-\xi_2$. The bracket in
\eqref{eq:alg2} is then evaluated as $-\mathrm{expm1}(\Delta\xi)\,\hat a_l\hat b_l$,
which is free of cancellation near chemical equilibrium. Otherwise all four
projections are computed.

\subsection{Limits and quadrature rules}

All rules are Gauss--Legendre (GL) with fixed numbers of nodes. Default
values are given in brackets.

\begin{description}
\item[$\sigma$ (i.e.\ $s$):] $\sigma\in[0,\sigma_{\max}]$, $\sigma_{\max}=\sqrt{z_{\rm cut}}$, where
  \begin{equation}
    z_{\rm cut} = n_{\rm cut}\,T_{\max}+\mu_{\max},\quad
    T_{\max}=\max_iT_i,\quad
    \mu_{\max}=\max\big(0,\;\xi_1T_1+\xi_2T_2,\;\xi_3T_3+\xi_4T_4\big),
  \end{equation}
  with $n_{\rm cut}=50$. This truncates $\sqrt s\le w_{\rm th}+z_{\rm cut}$,
  which is justified because $E\ge\sqrt s$ and $F\lesssim e^{-(E-\mu)/T}$.
  The range is split into geometrically graded panels
  \begin{equation}
    [0,\;\sigma_{\max}2^{1-K}],\quad [\sigma_{\max}2^{-j},\;\sigma_{\max}2^{1-j}]\;\;(j=K-1,\dots,1),
  \end{equation}
  with $K=\max\big(1,\lceil\log_2(\sigma_{\max}/\sigma_{\rm lo})\rceil\big)$ and
  $\sigma_{\rm lo}=\tfrac14\sqrt{\min(m_1,\dots,m_4,T_1,\dots,T_4)}$. This
  resolves both the threshold scale and the thermal scale. Each panel uses
  $n_s$ GL nodes [10].
\item[$\tau$ (i.e.\ $E$):] $\tau\in[0,\tau_{\max}]$, $\tau_{\max}=\sqrt{z_{\rm cut}}$,
  i.e.\ $\sqrt s\le E\le\sqrt s+z_{\rm cut}$. The range is split into two
  equal panels with $n_E$ GL nodes each [16].
\item[$c_1,\,c_3$:] $[-1,1]$ exactly. The same $n_c$ GL nodes $c_i$ with
  weights $w_i$ [32] are used for both; the projections are
  $\hat a_l=\sum_iw_iP_l(c_i)\,a(c_i)$ for $l=0,\dots,L$ [$L=24$].
\item[$\theta^*$:] $[0,\pi]$ exactly. The integrand is folded onto
  $[0,\pi/2]$ using $P_l(-x)=(-1)^lP_l(x)$:
  \begin{equation}
    M_l = \int_0^{\pi/2}\!\dd\theta\,\sin\theta\,P_l(\cos\theta)
    \big[\Msq(\theta^*{=}\theta)+(-1)^l\,\Msq(\theta^*{=}\pi-\theta)\big].
  \end{equation}
  The panels are graded towards $\theta=0$, where the $t$-channel peak sits
  (and, through the mirror term, the $u$-channel peak at $\theta^*=\pi$):
  \begin{equation}
    [0,\theta_p/8],\;[\theta_p/8,\theta_p/4],\;[\theta_p/4,\theta_p/2],\;\dots
    \;\text{(doubling, last panel ending at $\pi/2$)},
  \end{equation}
  with $\theta_p=\min\big(\pi/2,\;m_{\min}/\sqrt{k_1k_3}\big)$, the angular
  width of the exchange peak, and $m_{\min}=\max(\min_im_i,\,10^{-8}\sqrt s)$.
  Each panel uses $n_\theta$ GL nodes [16].
\item[$l$:] truncated at $L$.
\end{description}

\subsection{Accuracy}

Against the Maxwell--Boltzmann reference built from eq.~(A10), and against
runs with all node counts increased, the defaults give $\sim10^{-11}$
relative accuracy for $10^{-3}\le x\le50$ ($m_{A'}=2.5\,m_N$,
$\xi_{A'}=2\xi_N$). The exception is a boson close to condensation while
relativistic ($m_{A'}/T-\xi_{A'}\lesssim10^{-2}$, $T\gg m_{A'}$). There
$f_{A'}$ is sharply peaked in $c_1$ and the error is $\sim10^{-3}$;
$n_c=128$ reduces it to $\sim10^{-5}$.

%------------------------------------------------------------------------------
\section{Summary of integration variables}

\begin{center}
\small
\renewcommand{\arraystretch}{1.3}
\begin{tabular}{@{}l>{$}l<{$}>{$}l<{$}l@{}}
\toprule
 & \text{variable} & \text{range} & rule \\
\midrule
\multicolumn{4}{@{}l}{\emph{Algorithm 1} (outer $\to$ inner)}\\
 & E_1 & [m_1,\ \max(100T_1,100m_1)] & adaptive GK \\
 & E_2 & [\max(m_2,m_3{+}m_4{-}E_1),\ \max(100T_2,100m_2)] & adaptive GK \\
 & E_3 & [m_3,\ E_1{+}E_2{-}m_4] & adaptive GK \\
 & s   & [\max(s_{12}^-,s_{34}^-),\ \min(s_{12}^+,s_{34}^+)] & adaptive GK \\
 & t   & [t^-,\ t^+] & closed form \\
\midrule
\multicolumn{4}{@{}l}{\emph{Algorithm 2} (outer $\to$ inner)}\\
 & \sigma=\sqrt{\sqrt s-w_{\rm th}} & [0,\ \sqrt{z_{\rm cut}}] & GL, graded panels \\
 & \theta^* & [0,\ \pi] & GL, panels graded to $0,\pi$ \\
 & \tau=\sqrt{E-\sqrt s} & [0,\ \sqrt{z_{\rm cut}}] & GL, 2 panels \\
 & c_1,\ c_3 & [-1,\ 1] & GL \\
 & \phi & [0,\ \pi] & Funk--Hecke (exact), sum to $L$ \\
\bottomrule
\end{tabular}
\end{center}

\begin{thebibliography}{9}
\bibitem{Bringmann:2022aim}
T.~Bringmann, P.~F.~Depta, M.~Hufnagel, J.~Kersten, J.~T.~Ruderman and
K.~Schmidt-Hoberg, \emph{Minimal sterile neutrino dark matter},
Phys.\ Rev.\ D \textbf{107} (2023) L071702,
\href{https://arxiv.org/abs/2206.10630}{arXiv:2206.10630 [hep-ph]}.
\bibitem{Ala-Mattinen:2022nuj}
K.~Ala-Mattinen, M.~Heikinheimo, K.~Kainulainen and K.~Tuominen,
\emph{Momentum distributions of cosmic relics: Improved analysis},
Phys.\ Rev.\ D \textbf{105} (2022) 123005,
\href{https://arxiv.org/abs/2201.06456}{arXiv:2201.06456 [hep-ph]}.
\bibitem{Pons:1998}
J.~A.~Pons, J.~A.~Miralles and J.~M.~Ib\'a\~nez,
\emph{Legendre expansion of the neutrino-antineutrino annihilation kernel:
Influence of high order terms},
Astron.\ Astrophys.\ Suppl.\ Ser.\ \textbf{129} (1998) 343,
\href{https://arxiv.org/abs/astro-ph/9802333}{arXiv:astro-ph/9802333}.
\bibitem{Mueller:1998}
C.~M\"uller, \emph{Analysis of Spherical Symmetries in Euclidean Spaces},
Springer (1998).
\end{thebibliography}

\end{document}
